\begin{table*}[t]
\caption{Scope of related protection and attack studies.}
\label{tab:priorcomparison}\centering\footnotesize
\begin{tabularx}{\textwidth}{@{}>{\raggedright\arraybackslash}p{2.7cm}>{\raggedright\arraybackslash}X>{\raggedright\arraybackslash}X@{}}\toprule
Study & Protection or attack focus & Relationship to the present evidence\\\midrule
BioHashing~\cite{teoh2004biohashing} & Tokenized binary biometric transformation. & This paper separates hidden-operator linkage from compromised-operator inversion.\\[3pt]
IronMask~\cite{ironmask} & Angular face templates protected with real-valued coding. & Different construction; the binary-projection results in this paper do not evaluate its security.\\[3pt]
PolyProtect~\cite{polyprotect} & Polynomial mapping, including fully informed attack evaluation. & Both relate utility to protection; no common implementation or operating point was tested.\\[3pt]
MLP-Hash~\cite{mlphash} & Random multilayer mapping and binary output. & Nonlinear projection alternative; the attack in this paper is evaluated on three stated projection families.\\[3pt]
Deep-template benchmark~\cite{benchmark2026} & Recognition, unlinkability, and irreversibility across modalities. & Broader scheme coverage; this paper examines exposure count, length, and fresh-key digital acceptance.\\[3pt]
Wang et al.~\cite{wang2021csa} & Constraint-based pre-image attacks against protected face templates. & Related attack objective; the present study tests fresh-key acceptance after repeated exposure.\\[3pt]
Li and Hu~\cite{li2014multiplicity} & Recovery from multiple fingerprint records and known transformations. & Multiplicity precedent; the present face study adds utility, linkage, and conditional intervals.\\[3pt]
Mai / Vec2Face / adapter~\cite{mai2019reconstruction,vec2face,shahreza2024adapter} & Face-image reconstruction from embeddings using learned models. & Generated images and injected embeddings have different attack interfaces and denominators.\\[3pt]
Present study & RO, SRHT, PCA-guided RO; $K=1,2,5,10$; replay, linkage, inversion, and utility. & Joint storage, recognition, and repeated-exposure analysis for intended edge settings; no hardware benchmark.\\\bottomrule
\end{tabularx}
\tabnote{The comparison concerns study scope; external methods were not rerun under this protocol.}
\end{table*}
